\section{Residual and Horizon Tests}

Two controlled experiments were conducted to test whether the core problem---signal concentration in illiquid stocks---could be mitigated without introducing new data sources.

\subsection{Cross-Fitted Residual Target (RESID)}

The RESID experiment followed the methodology of Gu, Kelly, and Xiu (2020): instead of predicting raw forward returns, the model was trained to predict the residual from a cross-sectional regression of returns on four Chinese equity factors (market, size, value, turnover). The factors were constructed using cross-fitting to eliminate own-stock contamination.

\textbf{Results:} Residualization produced a modest improvement in the liquid U50 universe: U50 RESID net annual return improved from the RAW-H1 baseline. The low-liquidity concentration decreased directionally but remained substantial. The improvement was insufficient to cross the profitability gate: U50 RESID remained negative in absolute terms, and maximum drawdown remained above the 25\% threshold.

\textbf{Conclusion:} Residualization is retained as a cleaner research baseline but rejected as a strategy fix. The mild improvement confirms that some of the raw return variation is factor-driven, but removing this component does not substantially change the investability picture.

\subsection{Multi-Horizon Targets (H3 and H6)}

The multi-horizon experiment tested whether training on 3-month and 6-month cumulative returns could shift the signal into more liquid stocks. Leippold, Wang, and Zhou (2022) report that larger Chinese stocks show stronger predictability at longer horizons, motivating this test.

\textbf{Methodology:} Models were trained using 3-month ($R^{3m}_t = \prod_{k=1}^{3}(1+r_{t+k})-1$) and 6-month cumulative forward returns as targets. Training windows were purged (3 months for H3, 6 months for H6) to prevent overlap between training and validation return windows. The identical 20-seed ensemble, 3-fold structure, and executable portfolio engine were used.

\textbf{Important note:} An earlier version of the multi-horizon experiment (MH-1) contained a portfolio construction defect that produced invalid results with drawdowns exceeding 100\% (impossible for unlevered long-only portfolios). The defect was a capital-allocation error in the staggered cohort engine that caused each cohort to receive full capital rather than its proportional budget. The corrected experiment (MH-1R) used the proven monthly 90/80 rebalancing engine and produced the results reported here. The erroneous first results are preserved in the experiment ledger for methodological transparency.

\textbf{Results:} Neither horizon produced positive U50 net returns. H3 U50 returned \HorizonThreeLoss{} annually; H6 U50 returned \HorizonSixLoss{} annually. Turnover was lower at longer horizons (24--31\% for H6 vs.\ 53\% for H1), but the gross signal was insufficient regardless.

\textbf{Conclusion:} Multi-horizon targeting does not rescue OHLCV-only scalability. The liquidity-concentration problem is not a horizon artifact---it is a fundamental property of price-volume data in the tested sample.
